% Source: https://artofproblemsolving.com/community/c3h1644805_what_theorem_is_this
% Stable identity: aops-topic-1644805-q01
\textbf{The statement is false.} Take $a=2$ and $x=6$. Then $x-a=4$ does not divide $6$, but $a=2$ does divide $6$.

\textbf{What is true.} Since $x=(x-a)+a$, the number $x-a$ divides $x$ exactly when it divides $a$:
\[
(x-a)\mid x\iff(x-a)\mid a .
\]
So ``$x-a\nmid x$'' says that $x-a\nmid a$, which has nothing to do with whether $a\mid x$. (The claim is the contrapositive of ``if $a\mid x$ then $(x-a)\mid x$'', which fails for the same example.)

For real numbers, divisibility is not a useful notion at all: every nonzero real number $d$ ``divides'' $x$, since $\dfrac {x}{d}$ is a real number. That is why the statement is posed for integers.
